% Part: computability
% Chapter: computability-theory
% Section: applications-fixed-point

\documentclass[../../../include/open-logic-section]{subfiles}

\begin{document}

\olfileid{cmp}{thy}{apf}
\olsection{स्थिरबिंदू प्रमेयाचा उपयोग}

स्थिरबिंदू प्रमेयामुळे आंशिक संगणनक्षम फलनांची व्याख्या
त्यांच्या निर्देशांकांच्या साहाय्याने करता येते. उदाहरणार्थ,
प्रत्येक $y$ साठी पुढील अट पूर्ण करणारा निर्देशांक $e$
शोधता येतो:
\[
\cfind{e}(y) = e + y.
\]
आणखी एक उदाहरण म्हणजे स्थिरबिंदू प्रमेयाची सिद्धता
वापरून Java किंवा C++ मध्ये स्वतःचा मजकूर छापणारा
कार्यक्रम तयार करता येतो.

प्रत्येक $e$ साठी $W_e$ हे $\cfind{e}$ चे परिभाषाक्षेत्र
असू द्या. मग $W_0$, $W_1$, $W_2$,~\dots ही क्रमिका
संगणनक्षमपणे प्रगणनीय संचांचे प्रगणन करते, हे आठवा.
यांतील काही संच संगणनक्षम आहेत. आदान म्हणून $x$
घेऊन, $W_x$ संगणनक्षम असेल तेव्हा त्याच्या जातिबोधक
फलनाचा निर्देशांक परत करणारा अल्गोरिदम आहे का, असा
प्रश्न विचारता येतो. उत्तर ``नाही'' आहे; असा अल्गोरिदम नाही:

\begin{thm}
पुढील गुणधर्म असलेले कोणतेही आंशिक संगणनक्षम फलन $f$
नाही: $W_e$ संगणनक्षम असेल तेव्हा $f(e)$ परिभाषित
असते आणि $\cfind{f(e)}$ हे त्याचे जातिबोधक फलन असते.
\end{thm}

\begin{proof}
$f$ हे कोणतेही आंशिक संगणनक्षम फलन असू द्या. आपण असा
निर्देशांक $e$ तयार करू की $W_e$ संगणनक्षम असेल, पण
$\cfind{f(e)}$ हे त्याचे जातिबोधक फलन नसेल.
%READERNOTE{OLCMP-024 — गोठवलेल्या इंग्रजी प्रमेयात आंशिक संगणनक्षम f नसल्याचे म्हटले आहे, पण सिद्धतेची पहिली ओळ फक्त सर्वत्र परिभाषित संगणनक्षम f घेते. पुढील g ची रचना f(e) अपरिभाषित असतानाही आंशिक संगणनक्षम राहते; म्हणून येथे आवश्यक सर्वसाधारण आंशिक f घेतले आहे.}
स्थिरबिंदू प्रमेय वापरून असा निर्देशांक $e$ मिळतो की
\[
\cfind{e}(y) \simeq
\begin{cases}
0 & \text{जर $y=0$ आणि $\cfind{f(e)}(0) \fdefined = 0$ असेल तर} \\
\text{अपरिभाषित} & \text{अन्यथा.}
\end{cases}
\]
म्हणजे पुढीलप्रमाणे व्याख्यायित फलनावर स्थिरबिंदू प्रमेय
लावून $e$ मिळतो:
\[
g(x,y) \simeq
\begin{cases}
0 & \text{जर $y=0$ आणि $\cfind{f(x)}(0) \fdefined = 0$ असेल तर} \\
\text{अपरिभाषित} & \text{अन्यथा.}
\end{cases}
\]
$g$ आंशिक संगणनक्षम आहे हे अनौपचारिकपणे असे पाहता
येते: आदान $x$ आणि $y$ मिळाल्यावर अल्गोरिदम आधी
$y$ हे~$0$ शी समान आहे का ते तपासतो. समान असल्यास
तो $f(x)$ संगणित करतो आणि मग सार्वत्रिक यंत्र वापरून
$\cfind{f(x)}(0)$ संगणित करतो. हे शेवटचे संगणन थांबून
$0$ परत करत असेल, तर अल्गोरिदम~$0$ परत करतो;
अन्यथा अल्गोरिदम थांबत नाही.

आता $\cfind{f(e)}(0)$ परिभाषित असून $0$ च्या समान
असेल, तर $\cfind{e}(y)$ हे $y$~$0$ च्या समान असेल तेव्हाच
परिभाषित असते; म्हणून $W_e =
\{ 0 \}$. $\cfind{f(e)}(0)$ अपरिभाषित असेल किंवा
परिभाषित असूनही $0$ च्या समान नसेल, तर $W_e = \emptyset$.
दोन्ही स्थितींत अपेक्षित गुणधर्म फसतो: या आदानावर फलन
निर्देशांक परत करतच नसेल, तर अपेक्षित परिभाषितता मिळत
नाही; निर्देशांक परत करत असेल, तर $\cfind{f(e)}$ हे~$W_e$
चे जातिबोधक फलन नाही, कारण आदान $0$ वर ते चुकीचे
उत्तर देते.
\end{proof}

\end{document}
